\section{Related Work}
\label{sec:background:related}

\noindent\textit{This section surveys the prior research on
which the proposed solution rests. The grading rubric
(\cite{rubric}, Câu~3) requires at least three related problems
and at least three peer-reviewed studies; we identify four
distinct problem areas and seven supporting references, each of
which a later section of this chapter develops in technical
depth.}

\subsection{Bus operational analytics on GPS and AFC data}
\label{sec:background:related:ops}

The closest peer-reviewed work to the present thesis is
Srisakaokul~\textit{et~al.}'s study of public-city-bus service
quality in Bangkok~\cite{srisakaokul2023}, which derives three
quality-of-service tiers (QoS-1 trip completion, QoS-2 on-path
driving, QoS-3 on-schedule) from raw GPS pings and schedule
data. The framework transfers directly to the HCMC case and is
developed in Section~\ref{sec:theory:qos}.
A complementary line of research, exemplified by
Derevitskiy~\textit{et~al.}~\cite{derevitskiy2016}, treats the
bus fleet itself as a distributed traffic sensor: GPS pings are
aggregated into road-graph edge-speed estimates and propagated
to unsensed edges via macroscopic travel-time models. The
authors report a mean absolute percentage error around twenty
percent at distances under fifteen kilometres. The HCMC dataset's
fifty-day GPS window is in principle sufficient for the same
construction; the present thesis defers this extension to future
work (Section~\ref{sec:future:predictive}) in order to focus on
historical analytics in Phase~1.

A third line of work, Lin~\&~Zeng~\cite{linzeng2014}, builds a
real-time bus-arrival-time predictor on a link-node route
representation, with dwell at time-check stops as the dominant
predictor. The same time--distance trajectory construction
underlies the Marey diagram adopted in this thesis as the
Routes-card detail visualisation
(Section~\ref{sec:background:marey}).

\subsection{Visualisation of spatial--temporal traffic data}
\label{sec:background:related:viz}

The PhD work of Liu~\cite{liu2022} contributes three
visualisation primitives that the present thesis adopts
verbatim: \emph{Gridify} (a decoupling of data and visual grid
abstractions, used in the dashboard's hex-grid map),
\emph{GroupSet} (a set-intersection partitioning of routes by
archetype, used in the Routes card's UpSet matrix), and
\emph{FuzzyCut} (interactive categorisation). Liu also formalises
the Multiple Coordinated Views interaction model, which the
dashboard implements at full scope as the A1 linked-brushing
invariant of Chapter~\ref{chapter:proposed}.

A second body of work, the storytelling-with-data principles of
Knaflic~\cite{knaflic2015}, governs the chart-selection and
clutter-handling rules adopted in the dashboard. The A2 story
subtitle invariant of Chapter~\ref{chapter:proposed} is the
direct application of these principles as a design constraint
(no chart without a story subtitle; muted defaults; bright colour
reserved for the one element the story points at).

\subsection{Bus--metro intermodal connectivity and existing
HCMC tooling}
\label{sec:background:related:hcmc}

The advisor specification~\cite{thoai2025} names bus--metro
intermodal connectivity as a required feature of the dashboard
and supplies the table of nineteen bus routes that feed Metro
Line~1. The opening of Metro Line~1 on
22~December~2024~\cite{metro_wiki} sharpens the gap: the
effectiveness of urban metro mobility now depends on bus feeders
arriving on time at interchange stations, a connectivity
relationship the existing operational tools cannot monitor.

The principal existing HCMC analytical artefact is the MCPT
public route directory~\cite{ttgtcc2024}, which is a
navigational tool (routes can be looked up) but not an
analytical one (no aggregation, visualisation, or temporal
analysis). The HCMC dataset specification~\cite{hcmcgps2025}
documents the fifty-day GPS window of thirty-one routes and the
April~2025 parade week as a natural-experiment opportunity.
Together these references frame what the present thesis adds to
the local ecosystem: the first interactive analytical surface
over the network's GPS and AFC streams.
